\documentclass[aspectratio=169,11pt]{beamer}
\usepackage{fontspec}
\usetheme[progressbar=frametitle,numbering=fraction,sectionpage=none]{metropolis}
\IfFontExistsTF{Arial}{\setsansfont{Arial}\setmainfont{Arial}}{\setsansfont{TeX Gyre Heros}\setmainfont{TeX Gyre Heros}}
\usepackage{graphicx,booktabs,array,amsmath}
\definecolor{researchnavy}{HTML}{203A49}
\definecolor{researchteal}{HTML}{368C83}
\definecolor{researchpaper}{HTML}{FAFAF8}
\setbeamercolor{normal text}{fg=researchnavy,bg=researchpaper}
\setbeamercolor{frametitle}{bg=researchnavy,fg=white}
\setbeamercolor{alerted text}{fg=researchteal}
\setbeamercolor{progress bar}{fg=researchteal,bg=researchnavy!12}
\setbeamersize{text margin left=8mm,text margin right=8mm}
\setbeamerfont{frametitle}{size=\large,series=\bfseries}
\setbeamerfont{footnote}{size=\scriptsize}
\setbeamertemplate{navigation symbols}{}
\title{Water storage change across Thailand}
\subtitle{Climate reconstruction and the limits of human attribution}
\author{Farzeen Khan}
\institute{MS in Water Engineering and Management\\Asian Institute of Technology\\[2mm]Advisor: Dr. Natthachet Tangdamrongsub}
\date{October 2026}
\newcommand{\sourcenote}[1]{\vfill\par{\fontsize{7.1}{8.4}\selectfont\color{researchnavy!75}#1}}
\begin{document}

\begin{frame}[plain]\vspace{3mm}{\LARGE\bfseries Water storage change across Thailand}\par\vspace{3mm}{\large\color{researchteal}Climate reconstruction and the limits of human attribution}\par\vspace{7mm}{\large Farzeen Khan}\par\vspace{2mm}{\small MS in Water Engineering and Management\\Asian Institute of Technology}\par\vspace{5mm}{\small Advisor: Dr. Natthachet Tangdamrongsub}\par\vspace{2mm}{\small October 2026}\end{frame}
\begin{frame}{Research focus}
\begin{itemize}\setlength{\itemsep}{3mm}
\item How are total and groundwater storage changing across Thailand’s 22 principal basins?
\item How much of the observed change can the climate baseline explain?
\item How strongly do independent observations support the remaining residual?
\end{itemize}
\end{frame}
\begin{frame}{Research objectives}
\small\setlength{\tabcolsep}{3pt}\renewcommand{\arraystretch}{1.35}\begin{tabular}{>{\raggedright\arraybackslash}p{0.46\textwidth}>{\raggedright\arraybackslash}p{0.46\textwidth}}\toprule
\textbf{Objective} & \textbf{Output}\\\midrule
1  Storage change & Monthly TWSA/GWSA, Sen slopes and Mann–Kendall tests\\
2  Climate and residual & Climate reconstruction, GLDAS benchmark and observation checks\\
3  Decision support & Interactive basin app, input data and reproducible code\\
\bottomrule\end{tabular}
\end{frame}
\begin{frame}{Thailand’s 22 principal basins}
\begin{columns}[c]\begin{column}{.43\textwidth}\centering\includegraphics[height=.65\textheight]{figures/01_study_basins.pdf}\end{column}\begin{column}{.54\textwidth}\begin{itemize}
\item Official ONWR reporting framework
\item GRACE information extends across basin boundaries
\item All basin tables retain coverage and uncertainty
\end{itemize}\end{column}\end{columns}
\sourcenote{ONWR principal-basin geometry. Reporting units do not imply independent GRACE observations.}
\end{frame}
\begin{frame}{Data and scientific roles}
\small\setlength{\tabcolsep}{3pt}\renewcommand{\arraystretch}{1.35}\begin{tabular}{>{\raggedright\arraybackslash}p{0.46\textwidth}>{\raggedright\arraybackslash}p{0.46\textwidth}}\toprule
\textbf{Record} & \textbf{Role}\\\midrule
CSR RL06.3 GRACE/GRACE-FO & Observed total storage\\
GLDAS Noah and G3P & Non-groundwater storage compartments\\
Paired WaterGAP simulations & Naturalized and historical-human baselines\\
Native P/PET and CHIRPS & Weather reconstruction and rainfall sensitivity\\
Thai wells and MSEA-Res & Groundwater and reservoir consistency evidence\\
\bottomrule\end{tabular}
\sourcenote{Primary TWSA: April 2002–September 2024. Derived GWSA: April 2002–September 2023.}
\end{frame}
\begin{frame}{Analysis workflow}
\begin{itemize}\setlength{\itemsep}{3mm}
\item Harmonize monthly units, dates, area weights and the 2004–2009 reference
\item Calculate GWSA and estimate storage trends
\item Reconstruct naturalized storage using weather and climate memory
\item Calculate the residual, examine uncertainty and compare independent observations
\end{itemize}
\end{frame}
\begin{frame}{Groundwater storage balance}
\[ \mathrm{GWSA}=\mathrm{TWSA}-\mathrm{SMSA}-\mathrm{SWSA}-\mathrm{CWSA}-\mathrm{SWEA} \]
\begin{itemize}\setlength{\itemsep}{3mm}
\item Subtract soil, surface, canopy and snow storage anomalies
\item Use common valid cells and a common reference period
\item Groundwater uncertainty includes every subtracted compartment
\end{itemize}
\sourcenote{Storage anomalies use mm of equivalent water depth. Snow is small in Thailand but retained consistently.}
\end{frame}
\begin{frame}{Naturalized and historical simulations}
\small\setlength{\tabcolsep}{3pt}\renewcommand{\arraystretch}{1.35}\begin{tabular}{>{\raggedright\arraybackslash}p{0.46\textwidth}>{\raggedright\arraybackslash}p{0.46\textwidth}}\toprule
\textbf{Symbol} & \textbf{Meaning}\\\midrule
O & Observed or compartment-derived storage\\
C & Naturalized WaterGAP baseline\\
A & Historical WaterGAP with represented human activities\\
\bottomrule\end{tabular}
\sourcenote{Paired simulations share model structure and forcing. The naturalized baseline remains a model experiment.}
\end{frame}
\begin{frame}{Residual accounting}
\[ R=O-C=\underbrace{(A-C)}_{H}+\underbrace{(O-A)}_{E} \]
\begin{itemize}\setlength{\itemsep}{3mm}
\item H describes the human effect represented by the paired model
\item E retains the observation–historical-model discrepancy
\item The primary residual uses the published naturalized baseline
\end{itemize}
\sourcenote{The statistical emulator reconstructs C from weather. Its prediction does not replace C in the primary residual.}
\end{frame}
\begin{frame}{Chronological model evaluation}
\small\setlength{\tabcolsep}{3pt}\renewcommand{\arraystretch}{1.35}\begin{tabular}{>{\raggedright\arraybackslash}p{0.307\textwidth}>{\raggedright\arraybackslash}p{0.307\textwidth}>{\raggedright\arraybackslash}p{0.307\textwidth}}\toprule
\textbf{Period} & \textbf{Years} & \textbf{Use}\\\midrule
Warm-up & 1981–1990 & Initialize climate memory\\
Training & 1991–2009 & Fit transformations and tune models\\
Validation & 2010–2015 & Freeze model selection\\
Test & 2016–2023 & Evaluate unseen years\\
\bottomrule\end{tabular}
\sourcenote{Ridge, Elastic Net, splines and Random Forest. Predictors exclude GRACE and groundwater observations.}
\end{frame}
\begin{frame}{Precipitation and potential ET}
\centering\includegraphics[width=\textwidth,height=.64\textheight,keepaspectratio]{figures/forcing_national.pdf}
\sourcenote{National monthly values, 2003–2023. Native P/PET drive reconstruction. CHIRPS supplies the rainfall replacement test.}
\end{frame}
\begin{frame}{Spatial distribution of the meteorological inputs}
\centering\includegraphics[width=\textwidth,height=.64\textheight,keepaspectratio]{figures/forcing_spatial.pdf}
\sourcenote{Available-month basin means, 2003–2023. Values are mm/month, not annual totals. Each map uses its labelled scale.}
\end{frame}
\begin{frame}{Observed storage through time}
\centering\includegraphics[width=\textwidth,height=.64\textheight,keepaspectratio]{figures/storage_national.pdf}
\sourcenote{National CSR TWSA and compartment-derived GWSA. Missing months remain gaps. Groundwater ends earlier because of compartment availability.}
\end{frame}
\begin{frame}{National storage decline}
\small\setlength{\tabcolsep}{3pt}\renewcommand{\arraystretch}{1.35}\begin{tabular}{>{\raggedright\arraybackslash}p{0.307\textwidth}>{\raggedright\arraybackslash}p{0.307\textwidth}>{\raggedright\arraybackslash}p{0.307\textwidth}}\toprule
\textbf{Storage} & \textbf{Sen slope (mm/year)} & \textbf{Conditional 95\% interval}\\\midrule
Total storage & $-$7.51 & $-$9.79 to $-$5.68\\
Groundwater storage & $-$6.84 & $-$8.77 to $-$5.11\\
\bottomrule\end{tabular}
\sourcenote{237 TWSA months through September 2024; 225 GWSA months through September 2023. Full-record seasonal Sen slopes.}
\end{frame}
\begin{frame}{Storage trend patterns across basins}
\centering\includegraphics[width=\textwidth,height=.64\textheight,keepaspectratio]{figures/storage_slopes.pdf}
\sourcenote{Seasonal Sen slopes. Negative values indicate decline. All basins retain the CSR small-area caution.}
\end{frame}
\begin{frame}{Statistical evidence for basin trends}
\centering\includegraphics[width=\textwidth,height=.64\textheight,keepaspectratio]{figures/trend_evidence.pdf}
\sourcenote{Seasonal Mann–Kendall with basin FDR. Colours saturate at q = 0.10. This test does not remove serial dependence or product uncertainty.}
\end{frame}
\begin{frame}{Basin ranking depends on the processing chain}
\centering\includegraphics[width=\textwidth,height=.64\textheight,keepaspectratio]{figures/05_product_ranking_scatter.pdf}
\sourcenote{Common-date rank correlations: TWSA 0.138 and GWSA 0.027. Groundwater chains also differ in compartment processing.}
\end{frame}
\begin{frame}{Climate reconstruction performance}
\centering\includegraphics[width=\textwidth,height=.64\textheight,keepaspectratio]{figures/climate_scores.pdf}
\sourcenote{Ensemble test anomaly NSE against naturalized WaterGAP, 2016–2023. High scores verify reconstruction of the parent model.}
\end{frame}
\begin{frame}{Naturalized storage and weather reconstruction}
\centering\includegraphics[width=\textwidth,height=.64\textheight,keepaspectratio]{figures/climate_national.pdf}
\sourcenote{National series. Training: 1991–2009; validation: 2010–2015; test: 2016–2023. Parent-model agreement is distinct from observational validation.}
\end{frame}
\begin{frame}{Climate memory improves reconstruction}
\centering\includegraphics[width=\textwidth,height=.64\textheight,keepaspectratio]{figures/13_climate_memory_ablation.pdf}
\sourcenote{Controlled Ridge experiment. National test anomaly NSE: current weather / memory = 0.283 / 0.890 for TWSA and 0.039 / 0.723 for GWSA.}
\end{frame}
\begin{frame}{Performance varies between test years}
\centering\includegraphics[width=\textwidth,height=.64\textheight,keepaspectratio]{figures/annual_skill.pdf}
\sourcenote{2022 national anomaly NSE: TWSA 0.234; GWSA $-$1.067. Negative NSE indicates performance below the mean benchmark.}
\end{frame}
\begin{frame}{Residual and modeled human effects differ}
\centering\includegraphics[width=\textwidth,height=.64\textheight,keepaspectratio]{figures/residual_national.pdf}
\sourcenote{R = H + E on common observed months. The remaining discrepancy includes unresolved model and observation errors.}
\end{frame}
\begin{frame}{Spatial residual trends}
\centering\includegraphics[width=\textwidth,height=.64\textheight,keepaspectratio]{figures/residual_spatial.pdf}
\sourcenote{Calendar-month fixed-effect OLS slopes on matched support. These slopes describe residual change and do not measure pumping.}
\end{frame}
\begin{frame}{Additive national trend partition}
\small\setlength{\tabcolsep}{3pt}\renewcommand{\arraystretch}{1.35}\begin{tabular}{>{\raggedright\arraybackslash}p{0.307\textwidth}>{\raggedright\arraybackslash}p{0.307\textwidth}>{\raggedright\arraybackslash}p{0.307\textwidth}}\toprule
\textbf{Component} & \textbf{TWSA} & \textbf{GWSA}\\\midrule
Observed O & -8.36 & -7.28\\
Naturalized C & -0.54 & -0.21\\
Modeled human H & -0.49 & 0.06\\
Residual R & -7.81 & -7.07\\
Discrepancy E & -7.32 & -7.13\\
\bottomrule\end{tabular}
\sourcenote{OLS slopes, mm/year, using the same months. R = H + E. Seasonal Sen slopes are reported separately because their medians are not additive.}
\end{frame}
\begin{frame}{GLDAS comparison on matching compartments}
\centering\includegraphics[width=\textwidth,height=.64\textheight,keepaspectratio]{figures/gldas_national.pdf}
\sourcenote{Soil + canopy + snow only. National r = 0.969, NSE = 0.738, RMSE = 52.49 mm across 261 months. Groundwater and surface water are excluded.}
\end{frame}
\begin{frame}{GLDAS agreement varies among basins}
\centering\includegraphics[width=\textwidth,height=.64\textheight,keepaspectratio]{figures/gldas_scores.pdf}
\sourcenote{Matched land-storage benchmark. NSE colours saturate below $-$1 and RMSE above 100 mm. GLDAS also contributes to the groundwater subtraction.}
\end{frame}
\begin{frame}{Groundwater level and storage use different units}
\small\setlength{\tabcolsep}{3pt}\renewcommand{\arraystretch}{1.35}\begin{tabular}{>{\raggedright\arraybackslash}p{0.46\textwidth}>{\raggedright\arraybackslash}p{0.46\textwidth}}\toprule
\textbf{Observation} & \textbf{What the comparison establishes}\\\midrule
Groundwater head or depth (m) & Local aquifer response and direction\\
Regional storage anomaly (mm) & Water-mass change over a basin\\
Paired annual correlation & Agreement in timing on observed years\\
\bottomrule\end{tabular}
\sourcenote{A head-to-storage conversion requires aquifer storage coefficients and representative areas. Those inputs are unavailable here.}
\end{frame}
\begin{frame}{Well selection changes residual agreement}
\small\setlength{\tabcolsep}{3pt}\renewcommand{\arraystretch}{1.35}\begin{tabular}{>{\raggedright\arraybackslash}p{0.23\textwidth}>{\raggedright\arraybackslash}p{0.23\textwidth}>{\raggedright\arraybackslash}p{0.23\textwidth}>{\raggedright\arraybackslash}p{0.23\textwidth}}\toprule
\textbf{Station panel} & \textbf{TWS r} & \textbf{GWS r} & \textbf{GWS years}\\\midrule
Longest overlap & 0.127 & 0.037 & 17\\
Fixed eligibility / wider history & 0.657 & 0.511 & 17\\
Common period & 0.325 & 0.029 & 8\\
Strictly balanced & 0.305 & $\approx$0.000 & 8\\
\bottomrule\end{tabular}
\sourcenote{511 wells provide 5,484 aggregates. Co-located wells are grouped. The common-period panels cover eight usable storage years within 2013–2022.}
\end{frame}
\begin{frame}{Spatial groundwater consistency}
\centering\includegraphics[width=\textwidth,height=.64\textheight,keepaspectratio]{figures/head_spatial.pdf}
\sourcenote{Residual–head correlations on available annual pairs. Grey hatching marks insufficient evidence. These comparisons cannot validate a human-only total.}
\end{frame}
\begin{frame}{Independent reservoir evidence}
\centering\includegraphics[width=\textwidth,height=.64\textheight,keepaspectratio]{figures/reservoir_national.pdf}
\sourcenote{26-reservoir fixed panel, 73 scene-supported months. Adjusted r = 0.439, year-cluster 95\% interval 0.188–0.665. Both curves remove season and linear trend.}
\end{frame}
\begin{frame}{Reservoir agreement and observation support}
\centering\includegraphics[width=\textwidth,height=.64\textheight,keepaspectratio]{figures/reservoir_spatial.pdf}
\sourcenote{Scene-screened MSEA-Res Level 1. Reservoir storage checks a surface-water component. Climate and operating decisions both affect reservoirs.}
\end{frame}
\begin{frame}{Sensitivity to the rainfall source}
Test anomaly NSE in the frozen Ridge memory experiment\par\vspace{2mm}
\small\setlength{\tabcolsep}{3pt}\renewcommand{\arraystretch}{1.35}\begin{tabular}{>{\raggedright\arraybackslash}p{0.307\textwidth}>{\raggedright\arraybackslash}p{0.307\textwidth}>{\raggedright\arraybackslash}p{0.307\textwidth}}\toprule
\textbf{Storage} & \textbf{Native rainfall} & \textbf{CHIRPS replacement}\\\midrule
TWSA & 0.890 & 0.688\\
GWSA & 0.723 & 0.562\\
\bottomrule\end{tabular}
\sourcenote{PET and the native-forced parent target stay fixed. This tests forcing dependence, rather than which rainfall product is more accurate.}
\end{frame}
\begin{frame}{Decision support and reproducibility}
\begin{itemize}\setlength{\itemsep}{3mm}
\item Interactive storage histories, trend maps and observation support for all 22 basins
\item Downloadable basin performance tables and spatial/temporal forcing figures
\item Input subsets, pinned dependencies and executable scientific workflows
\end{itemize}
\sourcenote{https://pratikkhanal16.github.io/thailand-water-storage-dss/extended/  •  Source repository: pratikkhanal16/thailand-water-storage-dss}
\end{frame}
\begin{frame}{Conclusions}
\begin{itemize}\setlength{\itemsep}{3mm}
\item The primary records show national TWSA and GWSA decline
\item Basin rankings and residual interpretation depend on products and observation support
\item Climate memory reconstructs the naturalized parent well, with weaker performance in some years
\item Wells and reservoirs provide complementary, limited evidence for residual interpretation
\end{itemize}
\sourcenote{Contribution: a reproducible joint assessment of product sensitivity, climate memory, paired-model discrepancy and monitoring-network uncertainty.}
\end{frame}
\begin{frame}{Selected references}
\begin{itemize}\setlength{\itemsep}{3mm}
\item Save et al. (2016), CSR mascons. doi:10.1002/2016JB013007
\item Müller Schmied et al. (2024), WaterGAP. doi:10.5194/gmd-17-8817-2024
\item Sharifi et al. (2025), groundwater storage. doi:10.5194/hess-29-6985-2025
\item MSEA-Res (2025), reservoir storage. doi:10.5194/essd-17-2693-2025
\item GROW groundwater observations. doi:10.5281/zenodo.15149480
\end{itemize}
\sourcenote{Full bibliography, source terms and basin results accompany the report. Theme: Metropolis by Matthias Vogelgesang and contributors, https://github.com/matze/mtheme.}
\end{frame}
\appendix
\begin{frame}{Supporting results}
\begin{itemize}\setlength{\itemsep}{3mm}
\item Additional hydroclimatic context
\item Basin performance tables for questions
\item Observation support and interpretation
\end{itemize}
\end{frame}
\begin{frame}{Supporting hydroclimatic histories}
\centering\includegraphics[width=\textwidth,height=.64\textheight,keepaspectratio]{figures/supporting_temporal.pdf}
\sourcenote{National monthly observations, 2003–2023. Actual ET, soil moisture and surface temperature provide context; they are not added primary predictors.}
\end{frame}
\begin{frame}{Spatial hydroclimatic context}
\centering\includegraphics[width=\textwidth,height=.64\textheight,keepaspectratio]{figures/supporting_spatial.pdf}
\sourcenote{Basin available-month means, 2003–2023. Product coverage differs. Volumetric moisture is distinct from storage in mm.}
\end{frame}
\begin{frame}{Spatial support and uncertainty}
\begin{itemize}\setlength{\itemsep}{3mm}
\item National aggregation touches 68 native CSR mascons
\item Inverse weight concentration is about 50.1 and does not count independent observations
\item Every principal basin is below the provider’s approximate 200,000 km² small-area caution
\item Conditional intervals do not include all gravity, compartment and hydrological-model errors
\end{itemize}
\end{frame}
\begin{frame}{Observed storage trends for every basin (01–11)}
\fontsize{8.5}{10}\selectfont\setlength{\tabcolsep}{3pt}\renewcommand{\arraystretch}{1.0}\begin{tabular}{>{\raggedright\arraybackslash}p{0.126\textwidth}>{\raggedright\arraybackslash}p{0.126\textwidth}>{\raggedright\arraybackslash}p{0.126\textwidth}>{\raggedright\arraybackslash}p{0.126\textwidth}>{\raggedright\arraybackslash}p{0.126\textwidth}>{\raggedright\arraybackslash}p{0.126\textwidth}>{\raggedright\arraybackslash}p{0.126\textwidth}}\toprule
\textbf{Basin} & \textbf{TWS slope} & \textbf{TWS q} & \textbf{n} & \textbf{GWS slope} & \textbf{GWS q} & \textbf{n}\\\midrule
01 & -3.59 & <.001 & 237 & -2.88 & <.001 & 225\\
02 & -1.90 & 0.010 & 237 & -1.83 & 0.014 & 225\\
03 & -1.32 & 0.007 & 237 & -1.59 & <.001 & 225\\
04 & -3.65 & <.001 & 237 & -3.32 & <.001 & 225\\
05 & -4.42 & <.001 & 237 & -3.47 & <.001 & 225\\
06 & -3.71 & <.001 & 237 & -3.10 & <.001 & 225\\
07 & -3.21 & <.001 & 237 & -3.32 & <.001 & 225\\
08 & -3.79 & <.001 & 237 & -3.97 & <.001 & 225\\
09 & -3.74 & <.001 & 237 & -3.94 & <.001 & 225\\
10 & -7.36 & <.001 & 237 & -7.47 & <.001 & 225\\
11 & -5.50 & <.001 & 237 & -5.20 & <.001 & 225\\
\bottomrule\end{tabular}
\sourcenote{Slopes are seasonal Sen estimates in mm/year. q is seasonal Mann–Kendall FDR across basins; n is observed months. Full intervals and original p-values remain in CSR\_basin\_trends.csv. Basin names are listed in Appendix A.}
\end{frame}
\begin{frame}{Observed storage trends for every basin (12–22)}
\fontsize{8.5}{10}\selectfont\setlength{\tabcolsep}{3pt}\renewcommand{\arraystretch}{1.0}\begin{tabular}{>{\raggedright\arraybackslash}p{0.126\textwidth}>{\raggedright\arraybackslash}p{0.126\textwidth}>{\raggedright\arraybackslash}p{0.126\textwidth}>{\raggedright\arraybackslash}p{0.126\textwidth}>{\raggedright\arraybackslash}p{0.126\textwidth}>{\raggedright\arraybackslash}p{0.126\textwidth}>{\raggedright\arraybackslash}p{0.126\textwidth}}\toprule
\textbf{Basin} & \textbf{TWS slope} & \textbf{TWS q} & \textbf{n} & \textbf{GWS slope} & \textbf{GWS q} & \textbf{n}\\\midrule
12 & -7.06 & <.001 & 237 & -7.35 & <.001 & 225\\
13 & -6.89 & <.001 & 237 & -6.36 & <.001 & 225\\
14 & -5.25 & <.001 & 237 & -5.48 & <.001 & 225\\
15 & -8.64 & <.001 & 237 & -8.12 & <.001 & 225\\
16 & -7.68 & <.001 & 237 & -7.30 & <.001 & 225\\
17 & -7.51 & <.001 & 237 & -7.92 & <.001 & 225\\
18 & -5.85 & <.001 & 237 & -4.75 & <.001 & 225\\
19 & -34.37 & <.001 & 237 & -36.08 & <.001 & 225\\
20 & -6.25 & <.001 & 237 & -7.45 & <.001 & 225\\
21 & -1.62 & <.001 & 237 & -3.16 & <.001 & 225\\
22 & -21.48 & <.001 & 237 & -20.95 & <.001 & 225\\
\bottomrule\end{tabular}
\sourcenote{Slopes are seasonal Sen estimates in mm/year. q is seasonal Mann–Kendall FDR across basins; n is observed months. Full intervals and original p-values remain in CSR\_basin\_trends.csv. Basin names are listed in Appendix A.}
\end{frame}
\begin{frame}{TWSA climate reconstruction performance (01–11)}
\fontsize{8.5}{10}\selectfont\setlength{\tabcolsep}{3pt}\renewcommand{\arraystretch}{1.0}\begin{tabular}{>{\raggedright\arraybackslash}p{0.126\textwidth}>{\raggedright\arraybackslash}p{0.126\textwidth}>{\raggedright\arraybackslash}p{0.126\textwidth}>{\raggedright\arraybackslash}p{0.126\textwidth}>{\raggedright\arraybackslash}p{0.126\textwidth}>{\raggedright\arraybackslash}p{0.126\textwidth}>{\raggedright\arraybackslash}p{0.126\textwidth}}\toprule
\textbf{Basin} & \textbf{n} & \textbf{r} & \textbf{NSE} & \textbf{Anom NSE} & \textbf{RMSE} & \textbf{MAE}\\\midrule
01 & 96 & 0.991 & 0.981 & 0.807 & 22.08 & 16.90\\
02 & 96 & 0.990 & 0.979 & 0.904 & 17.32 & 13.88\\
03 & 96 & 0.990 & 0.973 & 0.812 & 27.37 & 19.39\\
04 & 96 & 0.981 & 0.961 & 0.821 & 15.61 & 11.58\\
05 & 96 & 0.974 & 0.946 & 0.732 & 21.43 & 15.40\\
06 & 96 & 0.991 & 0.981 & 0.876 & 15.90 & 12.00\\
07 & 96 & 0.987 & 0.974 & 0.868 & 15.24 & 11.08\\
08 & 96 & 0.987 & 0.972 & 0.873 & 15.38 & 11.43\\
09 & 96 & 0.992 & 0.983 & 0.901 & 14.14 & 10.83\\
10 & 96 & 0.967 & 0.928 & 0.792 & 42.23 & 30.13\\
11 & 96 & 0.990 & 0.980 & 0.875 & 15.34 & 12.27\\
\bottomrule\end{tabular}
\sourcenote{Frozen ensemble, 2016–2023 test data. Target: naturalized WaterGAP, not GRACE. Errors are mm; anomaly NSE removes the training monthly climatology. Higher NSE and lower errors are better.}
\end{frame}
\begin{frame}{TWSA climate reconstruction performance (12–22)}
\fontsize{8.5}{10}\selectfont\setlength{\tabcolsep}{3pt}\renewcommand{\arraystretch}{1.0}\begin{tabular}{>{\raggedright\arraybackslash}p{0.126\textwidth}>{\raggedright\arraybackslash}p{0.126\textwidth}>{\raggedright\arraybackslash}p{0.126\textwidth}>{\raggedright\arraybackslash}p{0.126\textwidth}>{\raggedright\arraybackslash}p{0.126\textwidth}>{\raggedright\arraybackslash}p{0.126\textwidth}>{\raggedright\arraybackslash}p{0.126\textwidth}}\toprule
\textbf{Basin} & \textbf{n} & \textbf{r} & \textbf{NSE} & \textbf{Anom NSE} & \textbf{RMSE} & \textbf{MAE}\\\midrule
12 & 96 & 0.987 & 0.974 & 0.864 & 11.22 & 8.35\\
13 & 96 & 0.974 & 0.942 & 0.778 & 32.00 & 23.99\\
14 & 96 & 0.991 & 0.981 & 0.856 & 18.49 & 14.13\\
15 & 96 & 0.975 & 0.942 & 0.764 & 28.65 & 22.72\\
16 & 96 & 0.979 & 0.957 & 0.705 & 14.60 & 11.45\\
17 & 96 & 0.987 & 0.972 & 0.742 & 16.75 & 13.62\\
18 & 96 & 0.979 & 0.957 & 0.827 & 14.38 & 10.97\\
19 & 96 & 0.982 & 0.965 & 0.911 & 17.38 & 13.55\\
20 & 96 & 0.982 & 0.963 & 0.903 & 22.54 & 18.01\\
21 & 96 & 0.977 & 0.951 & 0.887 & 25.44 & 18.76\\
22 & 96 & 0.985 & 0.970 & 0.903 & 17.36 & 13.16\\
\bottomrule\end{tabular}
\sourcenote{Frozen ensemble, 2016–2023 test data. Target: naturalized WaterGAP, not GRACE. Errors are mm; anomaly NSE removes the training monthly climatology. Higher NSE and lower errors are better.}
\end{frame}
\begin{frame}{GWSA climate reconstruction performance (01–11)}
\fontsize{8.5}{10}\selectfont\setlength{\tabcolsep}{3pt}\renewcommand{\arraystretch}{1.0}\begin{tabular}{>{\raggedright\arraybackslash}p{0.126\textwidth}>{\raggedright\arraybackslash}p{0.126\textwidth}>{\raggedright\arraybackslash}p{0.126\textwidth}>{\raggedright\arraybackslash}p{0.126\textwidth}>{\raggedright\arraybackslash}p{0.126\textwidth}>{\raggedright\arraybackslash}p{0.126\textwidth}>{\raggedright\arraybackslash}p{0.126\textwidth}}\toprule
\textbf{Basin} & \textbf{n} & \textbf{r} & \textbf{NSE} & \textbf{Anom NSE} & \textbf{RMSE} & \textbf{MAE}\\\midrule
01 & 96 & 0.982 & 0.963 & 0.830 & 5.74 & 4.47\\
02 & 96 & 0.986 & 0.956 & 0.901 & 4.98 & 4.08\\
03 & 96 & 0.982 & 0.953 & 0.781 & 10.59 & 7.70\\
04 & 96 & 0.969 & 0.933 & 0.755 & 7.04 & 5.33\\
05 & 96 & 0.963 & 0.927 & 0.678 & 8.87 & 6.32\\
06 & 96 & 0.980 & 0.958 & 0.842 & 3.94 & 2.92\\
07 & 96 & 0.977 & 0.948 & 0.824 & 4.83 & 3.96\\
08 & 96 & 0.984 & 0.963 & 0.865 & 5.52 & 4.22\\
09 & 96 & 0.982 & 0.958 & 0.832 & 6.04 & 4.78\\
10 & 96 & 0.968 & 0.935 & 0.661 & 4.90 & 4.09\\
11 & 96 & 0.958 & 0.914 & 0.627 & 7.09 & 5.40\\
\bottomrule\end{tabular}
\sourcenote{Frozen ensemble, 2016–2023 test data. Target: naturalized WaterGAP, not GRACE. Errors are mm; anomaly NSE removes the training monthly climatology. Higher NSE and lower errors are better.}
\end{frame}
\begin{frame}{GWSA climate reconstruction performance (12–22)}
\fontsize{8.5}{10}\selectfont\setlength{\tabcolsep}{3pt}\renewcommand{\arraystretch}{1.0}\begin{tabular}{>{\raggedright\arraybackslash}p{0.126\textwidth}>{\raggedright\arraybackslash}p{0.126\textwidth}>{\raggedright\arraybackslash}p{0.126\textwidth}>{\raggedright\arraybackslash}p{0.126\textwidth}>{\raggedright\arraybackslash}p{0.126\textwidth}>{\raggedright\arraybackslash}p{0.126\textwidth}>{\raggedright\arraybackslash}p{0.126\textwidth}}\toprule
\textbf{Basin} & \textbf{n} & \textbf{r} & \textbf{NSE} & \textbf{Anom NSE} & \textbf{RMSE} & \textbf{MAE}\\\midrule
12 & 96 & 0.981 & 0.962 & 0.812 & 4.62 & 3.61\\
13 & 96 & 0.959 & 0.920 & 0.703 & 5.50 & 4.45\\
14 & 96 & 0.985 & 0.970 & 0.812 & 4.39 & 3.12\\
15 & 96 & 0.963 & 0.902 & 0.360 & 10.51 & 8.46\\
16 & 96 & 0.970 & 0.933 & 0.564 & 10.88 & 8.87\\
17 & 96 & 0.980 & 0.952 & 0.589 & 11.67 & 9.60\\
18 & 96 & 0.958 & 0.907 & 0.711 & 6.84 & 5.11\\
19 & 96 & 0.968 & 0.926 & 0.823 & 10.40 & 8.46\\
20 & 96 & 0.946 & 0.892 & 0.668 & 11.69 & 9.70\\
21 & 96 & 0.953 & 0.904 & 0.808 & 9.19 & 7.33\\
22 & 96 & 0.967 & 0.922 & 0.780 & 11.32 & 8.94\\
\bottomrule\end{tabular}
\sourcenote{Frozen ensemble, 2016–2023 test data. Target: naturalized WaterGAP, not GRACE. Errors are mm; anomaly NSE removes the training monthly climatology. Higher NSE and lower errors are better.}
\end{frame}
\begin{frame}{GLDAS matched land storage performance (01–11)}
\fontsize{8.5}{10}\selectfont\setlength{\tabcolsep}{3pt}\renewcommand{\arraystretch}{1.0}\begin{tabular}{>{\raggedright\arraybackslash}p{0.126\textwidth}>{\raggedright\arraybackslash}p{0.126\textwidth}>{\raggedright\arraybackslash}p{0.126\textwidth}>{\raggedright\arraybackslash}p{0.126\textwidth}>{\raggedright\arraybackslash}p{0.126\textwidth}>{\raggedright\arraybackslash}p{0.126\textwidth}>{\raggedright\arraybackslash}p{0.126\textwidth}}\toprule
\textbf{Basin} & \textbf{n} & \textbf{r} & \textbf{NSE} & \textbf{RMSE} & \textbf{MAE} & \textbf{Bias}\\\midrule
01 & 261 & 0.949 & 0.850 & 55.42 & 46.06 & 4.24\\
02 & 261 & 0.955 & 0.877 & 38.82 & 31.16 & -3.20\\
03 & 261 & 0.957 & 0.670 & 62.72 & 55.91 & 3.44\\
04 & 261 & 0.940 & 0.702 & 56.19 & 48.08 & 11.12\\
05 & 261 & 0.931 & 0.620 & 65.72 & 57.34 & 12.09\\
06 & 261 & 0.961 & 0.901 & 37.15 & 30.56 & 3.75\\
07 & 261 & 0.962 & 0.874 & 36.09 & 29.37 & -2.62\\
08 & 261 & 0.926 & 0.712 & 59.38 & 51.34 & -2.20\\
09 & 261 & 0.956 & 0.825 & 48.01 & 41.12 & 1.04\\
10 & 261 & 0.871 & 0.390 & 77.84 & 69.70 & 7.74\\
11 & 261 & 0.900 & 0.756 & 59.88 & 49.64 & 3.91\\
\bottomrule\end{tabular}
\sourcenote{Soil, canopy and snow only; naturalized WaterGAP compared with GLDAS Noah on matched months. Errors are mm; bias = WaterGAP minus GLDAS. This does not validate full TWS, GWS or human-only storage.}
\end{frame}
\begin{frame}{GLDAS matched land storage performance (12–22)}
\fontsize{8.5}{10}\selectfont\setlength{\tabcolsep}{3pt}\renewcommand{\arraystretch}{1.0}\begin{tabular}{>{\raggedright\arraybackslash}p{0.126\textwidth}>{\raggedright\arraybackslash}p{0.126\textwidth}>{\raggedright\arraybackslash}p{0.126\textwidth}>{\raggedright\arraybackslash}p{0.126\textwidth}>{\raggedright\arraybackslash}p{0.126\textwidth}>{\raggedright\arraybackslash}p{0.126\textwidth}>{\raggedright\arraybackslash}p{0.126\textwidth}}\toprule
\textbf{Basin} & \textbf{n} & \textbf{r} & \textbf{NSE} & \textbf{RMSE} & \textbf{MAE} & \textbf{Bias}\\\midrule
12 & 261 & 0.848 & 0.465 & 75.38 & 66.62 & 7.09\\
13 & 261 & 0.814 & 0.317 & 95.06 & 85.91 & 13.97\\
14 & 261 & 0.955 & 0.788 & 62.51 & 53.84 & 4.43\\
15 & 261 & 0.934 & 0.621 & 60.08 & 51.92 & 5.74\\
16 & 261 & 0.860 & 0.453 & 78.69 & 68.66 & 9.02\\
17 & 261 & 0.891 & 0.559 & 71.30 & 62.24 & 3.90\\
18 & 261 & 0.858 & 0.514 & 80.68 & 71.40 & 13.50\\
19 & 261 & 0.866 & 0.690 & 49.28 & 39.89 & 0.01\\
20 & 261 & 0.844 & 0.544 & 53.29 & 41.99 & 10.73\\
21 & 261 & 0.678 & 0.403 & 62.21 & 51.62 & 9.14\\
22 & 261 & 0.887 & 0.686 & 49.60 & 40.29 & -6.04\\
\bottomrule\end{tabular}
\sourcenote{Soil, canopy and snow only; naturalized WaterGAP compared with GLDAS Noah on matched months. Errors are mm; bias = WaterGAP minus GLDAS. This does not validate full TWS, GWS or human-only storage.}
\end{frame}
\begin{frame}{Residual agreement with annual groundwater head (01–11)}
\fontsize{8.5}{10}\selectfont\setlength{\tabcolsep}{3pt}\renewcommand{\arraystretch}{1.0}\begin{tabular}{>{\raggedright\arraybackslash}p{0.126\textwidth}>{\raggedright\arraybackslash}p{0.126\textwidth}>{\raggedright\arraybackslash}p{0.126\textwidth}>{\raggedright\arraybackslash}p{0.126\textwidth}>{\raggedright\arraybackslash}p{0.126\textwidth}>{\raggedright\arraybackslash}p{0.126\textwidth}>{\raggedright\arraybackslash}p{0.126\textwidth}}\toprule
\textbf{Basin} & \textbf{TWS r} & \textbf{TWS detr r} & \textbf{Years} & \textbf{GWS r} & \textbf{GWS detr r} & \textbf{Years}\\\midrule
01 & 0.449 & 0.582 & 9 & -0.215 & 0.167 & 8\\
02 & 0.330 & 0.345 & 17 & 0.279 & 0.399 & 16\\
03 & -0.121 & -0.084 & 9 & 0.063 & -0.109 & 8\\
04 & -0.009 & -0.030 & 9 & -0.168 & -0.023 & 8\\
05 & 0.113 & 0.103 & 9 & -0.007 & 0.070 & 8\\
06 & 0.385 & 0.432 & 15 & 0.229 & 0.293 & 14\\
07 & -0.089 & -0.043 & 9 & -0.244 & 0.047 & 8\\
08 & 0.236 & 0.157 & 9 & 0.277 & -0.124 & 8\\
09 & 0.208 & 0.176 & 9 & 0.026 & -0.198 & 8\\
10 & 0.152 & 0.050 & 9 & 0.519 & 0.200 & 8\\
11 & — & — & 5 & — & — & 4\\
\bottomrule\end{tabular}
\sourcenote{Longest-overlap retained-flags panel. Positive r means lower residual tends to occur with lower head. Detr r removes a linear trend. An em dash means unavailable evidence, not zero. Head is not converted into storage; no cross-unit NSE or RMSE is reported. Other station panels and conditional intervals remain in the complete CSV.}
\end{frame}
\begin{frame}{Residual agreement with annual groundwater head (12–22)}
\fontsize{8.5}{10}\selectfont\setlength{\tabcolsep}{3pt}\renewcommand{\arraystretch}{1.0}\begin{tabular}{>{\raggedright\arraybackslash}p{0.126\textwidth}>{\raggedright\arraybackslash}p{0.126\textwidth}>{\raggedright\arraybackslash}p{0.126\textwidth}>{\raggedright\arraybackslash}p{0.126\textwidth}>{\raggedright\arraybackslash}p{0.126\textwidth}>{\raggedright\arraybackslash}p{0.126\textwidth}>{\raggedright\arraybackslash}p{0.126\textwidth}}\toprule
\textbf{Basin} & \textbf{TWS r} & \textbf{TWS detr r} & \textbf{Years} & \textbf{GWS r} & \textbf{GWS detr r} & \textbf{Years}\\\midrule
12 & — & — & 5 & — & — & 4\\
13 & -0.044 & -0.212 & 17 & 0.027 & -0.147 & 16\\
14 & 0.094 & 0.074 & 17 & 0.093 & 0.116 & 16\\
15 & 0.001 & -0.041 & 9 & -0.082 & -0.187 & 8\\
16 & — & — & 3 & — & — & 2\\
17 & 0.041 & 0.136 & 9 & -0.420 & -0.215 & 8\\
18 & -0.335 & -0.394 & 9 & -0.289 & -0.115 & 8\\
19 & 0.227 & 0.096 & 9 & 0.033 & -0.147 & 8\\
20 & -0.197 & -0.656 & 18 & -0.375 & -0.722 & 17\\
21 & — & — & 0 & — & — & 0\\
22 & -0.080 & -0.143 & 15 & -0.389 & -0.450 & 14\\
\bottomrule\end{tabular}
\sourcenote{Longest-overlap retained-flags panel. Positive r means lower residual tends to occur with lower head. Detr r removes a linear trend. An em dash means unavailable evidence, not zero. Head is not converted into storage; no cross-unit NSE or RMSE is reported. Other station panels and conditional intervals remain in the complete CSV.}
\end{frame}
\begin{frame}{Complete basin results and source data}
\begin{itemize}\setlength{\itemsep}{3mm}
\item Report Appendix C: 11 tables for every basin
\item Report Appendix D: 10 spatial and temporal forcing figures
\item The website supplies full-precision CSVs, vector figures and runnable code
\end{itemize}
\sourcenote{https://pratikkhanal16.github.io/thailand-water-storage-dss/extended/basin\_forcing/}
\end{frame}
\end{document}